% =============================================================================
% 1-pitch template — standalone-compilable lifecycle stage contract
%
% The pitch is venue-INDEPENDENT: the one-minute story before you pick a venue.
% RQs belong in 2-claims.tex (venue-coupled), not here.
%
% Layout follows ../../3-write-edit/_shared/sentence-format.md.
% Readability rules + section cues + worked examples: ref/pitch-readability.md.
% =============================================================================
\documentclass[11pt]{article}
\usepackage[margin=1in]{geometry}
\usepackage{parskip}
\usepackage{booktabs}
\usepackage{hyperref}
\title{Pitch: <Paper Title>}
\date{YYYY-MM-DD}
\begin{document}
\maketitle

\section*{Title}
% =========================================================
% Para [pitch.title] Working title
% Cue: short, specific, evocative; <=15 words; updated as the story sharpens.
% =========================================================
%% ---- P0.S1 ----
Working title for the paper.

\section*{One-Minute Pitch}
% =========================================================
% Para [pitch.kernel] One-minute pitch
% Cue: a short plain-language paragraph (~4-6 short sentences) for a NEWCOMER with no
%      background. Open with a framing sentence ("We study whether/how X relates to Y"),
%      then the puzzle, the method in plain words, the surprising finding, and why it
%      matters, so they understand it and feel interested. Not a single terse sentence.
% =========================================================
%% ---- P1.S1 ----
We study whether / how X relates to Y.
%
%% ---- P1.S2 ----
Plain-language context or the puzzle a newcomer needs.
%
%% ---- P1.S3 ----
The method in plain words.
%
%% ---- P1.S4 ----
The surprising finding.
%
%% ---- P1.S5 ----
Why it matters, and who could use it.

\section*{Hook}
% =========================================================
% Para [pitch.hook] Hook -- multiple candidate openings, all retained
% Cue: >=2 candidate methods (e.g. Paradox, Vivid Scene, Surprising Fact, Stakes, Gap).
%       Each candidate is a \subsection* block, one per narrative method.
%       ALL candidates are kept permanently displayed -- never collapsed.
%       Mark one as "(recommended lead)" but do NOT hide the rest.
%       The author chooses the final lead at write time.
% Hook section must show >=2 candidate methods, all retained, with a recommended lead marked.
% =========================================================

\subsection*{Candidate A: Paradox (recommended lead)}
%% ---- P2a.S1 ----
A paradox or tension that sets expectation against reality. 2-4 short sentences.

\subsection*{Candidate B: Vivid Scene}
%% ---- P2b.S1 ----
A concrete moment with specifics that makes the reader feel it. 2-4 short sentences.

\subsection*{Candidate C: Stakes}
%% ---- P2c.S1 ----
What is at risk -- opens with consequence. 2-4 short sentences.

% Add more candidates as needed. Each commits to ONE narrative move.
% See ref/pitch-readability.md for the full method menu.

\section*{Finding - Surprise}
% =========================================================
% Para [pitch.surprise] Surprise -- the non-obvious turn
% Cue: put the unexpected result in sentence 1, then the tension.
% =========================================================
%% ---- P3.S1 ----
The non-obvious turn, tension, or finding.

\section*{Implication - So What}
% =========================================================
% Para [pitch.implication] Implication -- so what, and who can use it
% Cue: name what changes and who can act in the first two sentences.
% =========================================================
%% ---- P4.S1 ----
What changes if this story is true, and who can use it.

\section*{Audience and Venue Fit}
% =========================================================
% Para [pitch.audience] Audience -- who the venue reaches and why they care
% Cue: name the reader and their need before the venue format.
% NOTE: this section is venue-AWARE (who reads that journal) but not
%       venue-COUPLED (RQs and claim framing live in 2-claims.tex).
% =========================================================
%% ---- P5.S1 ----
Who reads the target venue, and why this finding matters to that audience.

\section*{Evidence - Why Believe}
% =========================================================
% Para [pitch.evidence] Why believe -- evidence for each point
% Cue: tie each sentence to a table/display/model/check/source; mark planned as planned.
% =========================================================
%% ---- P6.S1 ----
Evidence for each point above, citing a source per claim (or intuition if seed-stage).

\section*{Limitation - Still Fragile}
% =========================================================
% Para [pitch.fragile] Still fragile -- the weakest point
% Cue: list only the top three highest-risk weaknesses.
% =========================================================
%% ---- P7.S1 ----
The weakest point or most important missing evidence.

\section*{Next Evidence Move}
% =========================================================
% Para [pitch.next] Next evidence move
% Cue: start with a verb and name the artifact this move updates.
% =========================================================
%% ---- P8.S1 ----
What discovery, task, probe, or review should happen next.

\end{document}
